\chapter{پیوست‌ها}
\label{ch:appendix}

در این بخش به مواردی می‌پرودازیم که دانستن آن‌ها در فرایند اپلای مهم بوده و با این حال در دسته‌بندی خاصی جا نگرفته یا در دسته‌بندی‌های زیادی
جا میگیرند. در صورت امکان محتوای این بخش به دسته‌بندی‌های مشخص خود منتقل خواهد شد تا از ایجاد شدن بخشی که محتوای نامربوط به هم دارد جلوگیری شود. 


\section{ترجمه}\label{sec:translation}
ارائه مدارک به مراجع رسمی خارج از کشور مثل دانشگاه هنگامی که این مدارک به زبانی غیر از زبان‌های مورد پذیرش 
کشور یا دانشگاه مقصد صادر شده‌اند نیاز به ترجمه دارد. معمولا ترجمه‌ها برای ارائه به مراجع رسمی باید بصورت \textbf{ترجمه رسمی} باشد
ولی موارد استثنایی هم وجود دارد. در ادامه انواع ترجمه توضیح داده می‌شود. 

\begin{itemize}
    \item \textbf{ترجمه رسمی}: این نوع ترجمه توسط دارالترجمه‌های رسمی و توسط مترجمان مورد تایید قوه‌قضاییه انجام می‌شود. 
    مدرک شما بعد از ترجمه روی کاغذ‌های مخصوصی و با مهر رسمی مترجم و بارکد بررسی اصالت ترجمه ارائه خواهد شد. برای انجام این نوع ترجمه باید
    مدرک شما حداقل‌هایی از رسمیت را داشته باشد. مثلا برای ترجمه رسمی ریزنمرات دانشگاهی باید ریزنمراتتان دارای مهر دانشگاه باشد. یا برای
    ترجمه رسمی فاکتور فروش یا خرید باید این فاکتور دارای مهر مغازه باشد.  

    \item \textbf{ترجمه غیر رسمی}: ترجمه غیررسمی دارای مهر و رسمیت خاصی نیست و می تواند توسط خود شما هم انجام شود. موارد معدودی
    وجود دارد که این نوع ترجمه در آن قابل استفاده است ولی تعداد آن‌ها صفر نیست.
\end{itemize}

مدرکی که آن را ترجمه رسمی می‌کنید بسته به جایی که قصد دارید آن را ارائه بدهید ممکن است نیاز به \textbf{مهرها و تاییدات} اضافه‌ای 
داشته باشد. در ادامه انواع ترجمه‌های تایید شده را توضیح می‌دهیم. 

\subsection{تاییدات ترجمه}\label{ssec:legalization}
مترجم با استفاده از مهر رسمی خود، صحت ترجمه خود را تایید می‌کند اما مراجع رسمی مخصوصا سفارت‌ها 
برای اطمینان از جعل نشدن مهر مترجم، صحت ترجمه و صحت سند از شما تاییدات دیگری علاوه بر مهر رسمی مترجم هم می‌خواهند. در ادامه انواع این تاییدات را 
توضیح می‌دهیم. 

\subsubsection{مهر دادگستری}
مهر قوه قضائیه یا مهر دادگستری، تأییدکننده هویت مترجم رسمی و اعتبار قانونی ترجمه است.
 این مهر نشان می‌دهد که ترجمه توسط مترجم مورد تأیید قوه قضائیه انجام شده و از نظر مراجع قانونی معتبر است ولی تایید کننده اصالت 
 محتوای ترجمه نیست.

\subsubsection{مهر وزارت امور خارجه}
این مهر تایید کننده اعتبار مدرک در کشور‌های خارجی است. معمولا مدارکی که قرار است به عنوان سند رسمی به سفارت ارائه شوند و در 
کشور مقصد اعتبار قانونی داشته باشند (مثلا مدرک دانشگاهی) باید این مهر را دریافت نمایند. توجه کنید که همه این مهر‌های 
ذکر شده فقط روی ترجمه زده می‌شوند و اصل مدرک تغییری نمی‌کند. 

\subsubsection{لگالیزیشن \lr{(Legalization)}}
این مهر که توسط سفارت کشور‌ها در ایران انجام می‌شود نشان‌ دهنده تایید شدن مدرک توسط سفارت است. این کار معمولا تنها برای مدارک خاصی انجام می‌شود
مثلا گواهی سوءپیشیته یا مدارک دانشگاهی و معمولا تنها پس از گرفتن تاییدات وزارتین (امور خارجه و دادگستری) قابل انجام است. 


برای گرفتن تاییدات معمولا مدرک باید قابل استعلام باشد یا مهرهای خاصی را دریافت کرده باشد که بسته به نوع مدرک فرق می‌کند. در ادامه توضیح می‌دهیم
هر مدرک برای گرفتن تاییدات وزارتین به مهر/استعلام چه مراجعی نیاز دارد. 

\begin{itemize}
    \item \textbf{دانشنامه/ریزنمرات}: مهر مترجم نیاز به مهر دانشگاه دارد ولی مهر وزارتین نیاز به تاییدیه ترجمه دارد که باید 
    از سامانه سجاد گرفته شده و پرینت آن کنار مدارک قرار بگیرد. 

    \item \textbf{شناسنامه}: وجود اصل مدرک برای همه مهر‌ها و تاییدات کافی است. 

    \item \textbf{گواهی عدم سوءپیشینه}: گواهی‌های سوءپیشینه که از سامانه عدل‌ایران گرفته می‌شوند نیاز به هیچ تایید اضافه‌ای ندارند 
    و می‌توانید تمام مهر‌های ترجمه را صرفا با اصل مدرک (که نیازی به مهر ندارد زیرا قابل استعلام است) بگیرید. 

    \item \textbf{فیش حقوقی/گواهی درآمدی}: این مدارک ابتدا باید مهر شرکتی که آن‌ها را صادر کرده دریافت کنند و برای دریافت 
    مهر وزارتین باید صفحه روزنامه رسمی مربوط به آن شرکت را از روزنامه رسمی کشور به صورت حضوری و با مهر برجسته دریافت کرده و به 
    ترجمه پیوست کنید. 

    \item \textbf{فاکتور خرید و فروش}: برای ترجمه رسمی این مدرک با مهر مترجم وجود مهر فروشگاه لازم و کافی است ولی برای
    گرفتن مهر وزارتین باید فاکتور مهر اتحادیه مربوطه را هم دریافت کند. البته ترجمه فاکتور معمولا به مهر وزارتین نیازی ندارد. 

\end{itemize}